\begin{proof}[Proof of \cref{obj:thm:fixed-interior-tightness-and-shrinking-radius-gap}]
\begin{enumerate}
\item Fix \(0<\epsilon<1/2\), and apply \cref{obj:thm:fixed-interior-tightness-and-shrinking-radius-gap-all-d} at this \(\epsilon\).  Choose the range constant
\[
c_\epsilon:=1 .
\]
Then \(c_\epsilon>0\).  The triangular benchmark algebra, the exact algebraic elbow, the residual wedge implication, and the diagonal witness assertions in the present statement are exactly the corresponding conclusions supplied by \cref{obj:thm:fixed-interior-tightness-and-shrinking-radius-gap-all-d}, with the same definitions of \(L_n\), \(b_{n,d}\), \(u_{n,d}\), \(q_{n,d,\sigma}\), \(\ell_{n,d,\sigma}\), and \(r_{n,d,\sigma}\).  In particular, the denominator in the residual-wedge ratio is the capped converse benchmark
\[
\frac1t+\min\left\{1,\frac{d_t}{t^2}\right\}
+\sigma_t^2\min\{1,u_{t,d_t}\},
\]
so the all-alphabet residual-wedge implication gives precisely
\[
\frac{b_{t,d_t}}{u_{t,d_t}}\to0,\qquad
u_{t,d_t}\to0,\qquad
\frac{b_{t,d_t}}{\sigma_t^2}\to0,\qquad
\sigma_t^2\to0 .
\]
For the witness \(d_n=n\), \(\sigma_n=L_n^{-1/2}\), the same all-alphabet conclusion gives membership in the residual wedge and the two displayed comparisons
\[
\frac{q_{n,n,L_n^{-1/2}}}{\ell_{n,n,L_n^{-1/2}}}\asymp L_n,
\qquad
\frac{q_{n,n,L_n^{-1/2}}}{\ell_{n,n,L_n^{-1/2}}}\to\infty .
\]
% lean: CausalSmith.Stat.DiscreteAteHeterogeneityFrontier.fixed_interior_tightness_and_shrinking_radius_gap_all_d

\item It remains only to specialize the fixed-interior and matching-elbow minimax clauses.  Let \(\underline\sigma\in(0,2]\).  The fixed-radius clause of \cref{obj:thm:fixed-interior-tightness-and-shrinking-radius-gap-all-d} supplies constants \(c,C\) with \(0<c\le C\) such that, for all \(n,d\in\mathbb N\), \(M,\sigma\in\mathbb R\),
\[
n>0,\qquad d>0,\qquad M\ge1,\qquad
0\le\sigma,\qquad \underline\sigma\le\sigma,\qquad \sigma\le2
\]
imply
\[
cM^2r_{n,d,\sigma}
\le
\mathsf R_{n,d,\epsilon,M,\sigma}
\le
CM^2r_{n,d,\sigma}.
\]
Since the implication holds for every positive \(d\), it also holds under the additional displayed range condition
\[
d\le c_\epsilon\frac{n^2}{L_n}
=\frac{n^2}{L_n}.
\]
Thus the fixed-interior-radii conclusion follows with the same constants \(c,C\).  % lean: CausalSmith.Stat.DiscreteAteHeterogeneityFrontier.fixed_interior_tightness_and_shrinking_radius_gap

\item Now fix \(K\ge0\).  The matching-elbow clause of \cref{obj:thm:fixed-interior-tightness-and-shrinking-radius-gap-all-d} supplies constants \(c,C\) with \(0<c\le C\) such that, whenever
\[
n>0,\qquad d>0,\qquad M\ge1,\qquad 0\le\sigma\le2
\]
and at least one of
\[
u_{n,d}\ge1,\qquad
u_{n,d}\le Kb_{n,d},\qquad
\sigma^2\le Kb_{n,d}
\]
holds, one has
\[
cM^2r_{n,d,\sigma}
\le
\mathsf R_{n,d,\epsilon,M,\sigma}
\le
CM^2r_{n,d,\sigma}.
\]
This is exactly the matching-elbows assertion in the present theorem.  Combining this with the clauses recorded in the first step and the choice \(c_\epsilon=1\) proves all asserted conclusions.  % lean: CausalSmith.Stat.DiscreteAteHeterogeneityFrontier.fixed_interior_tightness_and_shrinking_radius_gap
\end{enumerate}
\end{proof}
